% Exercise 3.15
\begin{itemize}
    \item[(a)] We have the following summary details for the least squares regression
        for the predictor \verb|zn|:
    \begin{lstlisting}
Residual standard error: 8.435 on 504 degrees of freedom
Multiple R-squared:  0.04019,   Adjusted R-squared:  0.03828
F-statistic:  21.1 on 1 and 504 DF,  p-value: 5.506e-06
    \end{lstlisting}

        For the predictor \verb|indus|:
    \begin{lstlisting}
Residual standard error: 7.866 on 504 degrees of freedom
Multiple R-squared:  0.1653,    Adjusted R-squared:  0.1637
F-statistic: 99.82 on 1 and 504 DF,  p-value: < 2.2e-16
    \end{lstlisting}

        For the predictor \verb|chas|:
    \begin{lstlisting}
Residual standard error: 8.597 on 504 degrees of freedom
Multiple R-squared:  0.003124,  Adjusted R-squared:  0.001146
F-statistic: 1.579 on 1 and 504 DF,  p-value: 0.2094
    \end{lstlisting}

        For the predictor \verb|nox|:
    \begin{lstlisting}
Residual standard error: 7.81 on 504 degrees of freedom
Multiple R-squared:  0.1772,    Adjusted R-squared:  0.1756
F-statistic: 108.6 on 1 and 504 DF,  p-value: < 2.2e-16
    \end{lstlisting}

        For the predictor \verb|rm|:
    \begin{lstlisting}
Residual standard error: 8.401 on 504 degrees of freedom
Multiple R-squared:  0.04807,   Adjusted R-squared:  0.04618
F-statistic: 25.45 on 1 and 504 DF,  p-value: 6.347e-07
    \end{lstlisting}

        For the predictor \verb|age|:
    \begin{lstlisting}
Residual standard error: 8.057 on 504 degrees of freedom
Multiple R-squared:  0.1244,    Adjusted R-squared:  0.1227
F-statistic: 71.62 on 1 and 504 DF,  p-value: 2.855e-16
    \end{lstlisting}

        For the predictor \verb|dis|:
    \begin{lstlisting}
Residual standard error: 7.965 on 504 degrees of freedom
Multiple R-squared:  0.1441,    Adjusted R-squared:  0.1425
F-statistic: 84.89 on 1 and 504 DF,  p-value: < 2.2e-16
    \end{lstlisting}

        For the predictor \verb|rad|:
    \begin{lstlisting}
Residual standard error: 6.718 on 504 degrees of freedom
Multiple R-squared:  0.3913,    Adjusted R-squared:   0.39
F-statistic: 323.9 on 1 and 504 DF,  p-value: < 2.2e-16
    \end{lstlisting}

        For the predictor \verb|tax|:
    \begin{lstlisting}
Residual standard error: 6.997 on 504 degrees of freedom
Multiple R-squared:  0.3396,    Adjusted R-squared:  0.3383
F-statistic: 259.2 on 1 and 504 DF,  p-value: < 2.2e-16
    \end{lstlisting}

        For the predictor \verb|ptratio|:
    \begin{lstlisting}
Residual standard error: 8.24 on 504 degrees of freedom
Multiple R-squared:  0.08407,   Adjusted R-squared:  0.08225
F-statistic: 46.26 on 1 and 504 DF,  p-value: 2.943e-11
    \end{lstlisting}

        For the predictor \verb|lstat|:
    \begin{lstlisting}
Residual standard error: 7.664 on 504 degrees of freedom
Multiple R-squared:  0.2076,    Adjusted R-squared:  0.206
F-statistic:   132 on 1 and 504 DF,  p-value: < 2.2e-16
    \end{lstlisting}

        For the predictor \verb|medv|:
    \begin{lstlisting}
Residual standard error: 7.934 on 504 degrees of freedom
Multiple R-squared:  0.1508,    Adjusted R-squared:  0.1491
F-statistic: 89.49 on 1 and 504 DF,  p-value: < 2.2e-16
        \end{lstlisting}

        We can safely rule out the effects of \verb|chas|, since the $p$-value is
        given by 0.2094, while the most significant association seems to be from
        \verb|dis|, \verb|rad|, and \verb|medv|, all with 
        $p < 2.2 \times {10}^{-16}$ and relatively low $F$-statistics.

        In Figure~\ref{fig3_15plot}, we can see these association clearly.
        \begin{figure}[!ht]
            \includegraphics[width=\textwidth, center]{../plots/ex3_15_a.pdf}
            \caption{Plots for models from  3.15.(a) \label{fig3_15plot}}
        \end{figure}
    \item[(b)] We have the following least squares regression for predicting
        \verb|crim| using all of the predictors:
    \begin{lstlisting}
> lm.fit <- lm(crim ~ ., data = Boston)
> summary(lm.fit)

Call:
lm(formula = crim ~ ., data = Boston)

Residuals:
     Min     1Q Median     3Q    Max
-8.534 -2.248 -0.348  1.087 73.923

Coefficients:
                            Estimate Std. Error t value Pr(>|t|)
(Intercept) 13.7783938  7.0818258   1.946 0.052271 .
zn           0.0457100  0.0187903   2.433 0.015344 *
indus       -0.0583501  0.0836351  -0.698 0.485709
chas        -0.8253776  1.1833963  -0.697 0.485841
nox         -9.9575865  5.2898242  -1.882 0.060370 .
rm           0.6289107  0.6070924   1.036 0.300738
age         -0.0008483  0.0179482  -0.047 0.962323
dis         -1.0122467  0.2824676  -3.584 0.000373 ***
rad          0.6124653  0.0875358   6.997 8.59e-12 ***
tax         -0.0037756  0.0051723  -0.730 0.465757
ptratio     -0.3040728  0.1863598  -1.632 0.103393
lstat        0.1388006  0.0757213   1.833 0.067398 .
medv        -0.2200564  0.0598240  -3.678 0.000261 ***
---
Signif. codes:  0 ‘***’ 0.001 ‘**’ 0.01 ‘*’ 0.05 ‘.’ 0.1 ‘ ’ 1

Residual standard error: 6.46 on 493 degrees of freedom
Multiple R-squared:  0.4493,    Adjusted R-squared:  0.4359
F-statistic: 33.52 on 12 and 493 DF,  p-value: < 2.2e-16
    \end{lstlisting}
        The multiple regression has an $F$-statistic of 33.52 on 12 and 493
        degrees of freedom, with a $p$-value of $p < 2.2 \times {10}^{-16}$, so
        that we have strong evidence for rejecting the null hypothesis
        $H_0 : \beta_1 = \beta_2 = \cdots = \beta_{12} = 0$, i.e., at least one
        of the predictors is associated with \verb|crim|. Looking at the
        individual predictors, we can reject the null hypothesis
        $H_0 : \beta_j = 0$ at the 1\% level for the predictors \verb|dis|
        ($p = 0.000373$), \verb|rad| ($p = 8.59 \times {10}^{-12}$), and
        \verb|medv| ($p = 0.000261$). The predictor \verb|zn| is also
        significant at the 5\% level ($p = 0.0153$), while none of the remaining
        predictors are significant at the 5\% level.
    \item[(c)] The results from (a) and (b) both agree on the importance of the
        predictors \verb|dis|, \verb|rad|, and \verb|medv| in predicting per
        capita crime rate. The fact that the multiple linear regression and the
        univariate regressions agree on the predictors provides even strong
        evidence for rejecting the null hypothesis for these predictors.

        In Figure~\ref{fig3_15coefs}, we plot the univariate regression
        coefficients from (a) against the multiple regression coefficients
        from (b).
        \begin{figure}[!ht]
            \includegraphics[scale=0.6, center]{../plots/ex3_15_c.pdf}
            \caption{Univariate vs multiple regression coefficients for 3.15.(c)
            \label{fig3_15coefs}}
        \end{figure}
    \item[(d)] We have the following coefficients for the non-linear least
        squares regression \verb|lm(crim ~ poly(zn, 3))| for the predictor
        \verb|zn|:
    \begin{lstlisting}
Coefficients:
                         Estimate Std. Error t value Pr(>|t|)
(Intercept)    3.6135     0.3722   9.709  < 2e-16 ***
poly(zn, 3)1 -38.7498     8.3722  -4.628  4.7e-06 ***
poly(zn, 3)2  23.9398     8.3722   2.859  0.00442 **
poly(zn, 3)3 -10.0719     8.3722  -1.203  0.22954
    \end{lstlisting}

        For the predictor \verb|indus|:
    \begin{lstlisting}
Coefficients:
                                Estimate Std. Error t value Pr(>|t|)
(Intercept)        3.614      0.330  10.950  < 2e-16 ***
poly(indus, 3)1   78.591      7.423  10.587  < 2e-16 ***
poly(indus, 3)2  -24.395      7.423  -3.286  0.00109 **
poly(indus, 3)3  -54.130      7.423  -7.292  1.2e-12 ***
    \end{lstlisting}

        For the predictor \verb|nox|:
    \begin{lstlisting}
Coefficients:
                            Estimate Std. Error t value Pr(>|t|)
(Intercept)     3.6135     0.3216  11.237  < 2e-16 ***
poly(nox, 3)1  81.3720     7.2336  11.249  < 2e-16 ***
poly(nox, 3)2 -28.8286     7.2336  -3.985 7.74e-05 ***
poly(nox, 3)3 -60.3619     7.2336  -8.345 6.96e-16 ***
    \end{lstlisting}

        For the predictor \verb|rm|:
    \begin{lstlisting}
Coefficients:
                         Estimate Std. Error t value Pr(>|t|)
(Intercept)    3.6135     0.3703   9.758  < 2e-16 ***
poly(rm, 3)1 -42.3794     8.3297  -5.088 5.13e-07 ***
poly(rm, 3)2  26.5768     8.3297   3.191  0.00151 **
poly(rm, 3)3  -5.5103     8.3297  -0.662  0.50858
    \end{lstlisting}

        For the predictor \verb|age|:
    \begin{lstlisting}
Coefficients:
                            Estimate Std. Error t value Pr(>|t|)
(Intercept)     3.6135     0.3485  10.368  < 2e-16 ***
poly(age, 3)1  68.1820     7.8397   8.697  < 2e-16 ***
poly(age, 3)2  37.4845     7.8397   4.781 2.29e-06 ***
poly(age, 3)3  21.3532     7.8397   2.724  0.00668 **
    \end{lstlisting}

        For the predictor \verb|dis|:
    \begin{lstlisting}
Coefficients:
                            Estimate Std. Error t value Pr(>|t|)
(Intercept)     3.6135     0.3259  11.087  < 2e-16 ***
poly(dis, 3)1 -73.3886     7.3315 -10.010  < 2e-16 ***
poly(dis, 3)2  56.3730     7.3315   7.689 7.87e-14 ***
poly(dis, 3)3 -42.6219     7.3315  -5.814 1.09e-08 ***
    \end{lstlisting}

        For the predictor \verb|rad|:
    \begin{lstlisting}
Coefficients:
                            Estimate Std. Error t value Pr(>|t|)
(Intercept)     3.6135     0.2971  12.164  < 2e-16 ***
poly(rad, 3)1 120.9074     6.6824  18.093  < 2e-16 ***
poly(rad, 3)2  17.4923     6.6824   2.618  0.00912 **
poly(rad, 3)3   4.6985     6.6824   0.703  0.48231
    \end{lstlisting}

        For the predictor \verb|tax|:
    \begin{lstlisting}
Coefficients:
                            Estimate Std. Error t value Pr(>|t|)
(Intercept)     3.6135     0.3047  11.860  < 2e-16 ***
poly(tax, 3)1 112.6458     6.8537  16.436  < 2e-16 ***
poly(tax, 3)2  32.0873     6.8537   4.682 3.67e-06 ***
poly(tax, 3)3  -7.9968     6.8537  -1.167    0.244
    \end{lstlisting}

        For the predictor \verb|ptratio|:
    \begin{lstlisting}
Coefficients:
                                    Estimate Std. Error t value Pr(>|t|)
(Intercept)          3.614      0.361  10.008  < 2e-16 ***
poly(ptratio, 3)1   56.045      8.122   6.901 1.57e-11 ***
poly(ptratio, 3)2   24.775      8.122   3.050  0.00241 **
poly(ptratio, 3)3  -22.280      8.122  -2.743  0.00630 **
    \end{lstlisting}

        For the predictor \verb|lstat|:
    \begin{lstlisting}
Coefficients:
                                Estimate Std. Error t value Pr(>|t|)
(Intercept)       3.6135     0.3392  10.654   <2e-16 ***
poly(lstat, 3)1  88.0697     7.6294  11.543   <2e-16 ***
poly(lstat, 3)2  15.8882     7.6294   2.082   0.0378 *
poly(lstat, 3)3 -11.5740     7.6294  -1.517   0.1299
    \end{lstlisting}

        For the predictor \verb|medv|:
    \begin{lstlisting}
Coefficients:
                             Estimate Std. Error t value Pr(>|t|)
(Intercept)       3.614      0.292  12.374  < 2e-16 ***
poly(medv, 3)1  -75.058      6.569 -11.426  < 2e-16 ***
poly(medv, 3)2   88.086      6.569  13.409  < 2e-16 ***
poly(medv, 3)3  -48.033      6.569  -7.312 1.05e-12 ***
    \end{lstlisting}

        A third order polynomial cannot be used for the predictor \verb|chas|
        with only two possible values (1 = next to the Charles river, 0 = not).
        In all other cases, the polynomial terms lead to an improvement in the
        model fit, although in some cases (\verb|lstat|, \verb|tax|,
        \verb|rad|, \verb|rm|, \verb|zn|), the third order term have
        significant $p$-values, indicating that a quadratic polynomial would
        provide a better fit for these cases.
\end{itemize}
